\subsection{From Frobenius slopes to Hodge types}\label{subsec:slopes-to-types}

We now combine the Frobenius polynomial with the theta parameters to
determine the Hodge types of the relevant Hecke spaces. The polynomial
will force a single Newton slope on each full generalized character
space. To turn this into a Hodge type, we also need its projector to
preserve the Hodge filtration. We first construct these projectors and
explain the filtered argument.

Let \(\mathbb T\) be the image of the full spherical Hecke algebra in
\(\End_{\Q}H^1(A_S,\Q)\). It is a finite-dimensional commutative
\(\Q\)-algebra; its active subalgebra is the image of
\(\mathcal H_w^0\). A numerical character is a \(\Q\)-algebra
homomorphism \(\chi:\mathbb T\to\Qbar\). Inside a character value,
the symbol of an endomorphism denotes its degree-one pullback.
After a finite extension of algebraic
coefficients, \(\mathbb T\) is a product of local Artinian algebras.
Its product idempotents \(p_\chi\) give the simultaneous generalized
character spaces
\[
H_\chi=p_\chi H^1(A_S,\Qbar).
\]
If \(\eta\) is a character of the active subalgebra, its generalized
space is
\begin{equation}\label{eq:active-generalized-space}
H^{\mathrm{act}}_\eta
=\bigoplus_{\chi|_{\mathrm{act}}=\eta}H_\chi.
\end{equation}
Indeed the image of the active idempotent in a local factor of the full
algebra is either zero or one, as decided by its residue character.
Thus the sum includes every full extension of \(\eta\) and all of its
nilpotent parts. No semisimplicity of the Hecke operators is assumed.

Each Hecke operator is induced by a rational endomorphism of \(A_S\)
defined over \(K\). In degree-one Betti, de Rham, and crystalline
cohomology, pullback is additive in abelian homomorphisms: in the
Künneth decomposition of \(H^1(A\times A)\), pullback by addition has
both components equal to the identity, as restriction to the two axes
shows. Pullback reverses composition, but the Hecke endomorphisms
commute. Pullback on Betti \(H^1\) is faithful by
Lemma~\ref{lem:weight-one-maps}, so polynomial relations in
\(\mathbb T\) are also relations among the corresponding
endomorphisms. The same degree-one functoriality preserves those
relations in de Rham and crystalline cohomology. Thus the \(p_\chi\)
remain idempotents in these degree-one
realizations after extension of coefficients. They need not have
rational coefficients; what matters is that they are
algebraic-coefficient polynomials in extending endomorphisms.

The following lemma applies to precisely these projected summands. Its
use of both complementary projectors is what turns weak admissibility
into equality of the Hodge and Newton numbers on each summand.

\begin{lemma}\label{lem:filtered-hecke-summands}
Let \(\mathcal A/\Z_q\) be an abelian scheme with generic fiber \(A\), and
put
\[
D_A=H^1_{\mathrm{dR}}(A/\Q_q).
\]
Equip \(D_A\) with its Hodge filtration and its good-reduction crystalline
Frobenius \(\varphi_A\). Then \(D_A\) is weakly admissible, its filtration
has only indices \(0\) and \(1\), and endomorphisms extending over \(\Z_q\)
preserve the filtration and commute with \(\varphi_A\).

Let \(L/\Q_q\) be finite. Suppose that an \(L\)-linear idempotent
\(p\) on \(D_A\otimes_{\Q_q}L\) preserves the filtration and commutes
with \(\varphi_A\otimes1\). If every Frobenius slope on its image is
\(t\in\{0,1\}\), then the induced filtration satisfies
\[
F^1\bigl(p(D_A\otimes L)\bigr)=
\begin{cases}
0,&t=0,\\
p(D_A\otimes L),&t=1.
\end{cases}
\]
\end{lemma}

\begin{proof}
Good-reduction crystalline comparison identifies
\[
D_{\mathrm{cris}}\!
 \left(H^1_{\mathrm{\acute et}}(A_{\overline{\Q}_q},\Q_q)\right)
\simeq
H^1_{\mathrm{cris}}(\mathcal A_{\Fq}/\Z_q)[1/q]
\]
and identifies the filtered de Rham realization with \(D_A\).
These identifications are functorial for extending endomorphisms and
compatible with Frobenius and filtration
\cite[Theorem~1.1(i), Theorem~12.1, and \S12.2]{BhattMorrowScholze2018}.\footnote{For the proper formal model and the algebraic/analytic
cohomology identifications, see \cite[\S7.2, p.~161]{Bosch2014FormalRigid}
and \cite[Theorems~9.1 and~9.3]{Scholze2013RigidAnalyticHodge}, read with
\cite{Scholze2016RigidAnalyticHodgeCorrigendum}. Passage from the algebraic
closure to its completion uses \cite[Tag~0DDG]{StacksProject}.
The \'etale comparisons are applied with \(\Z/q^n\)-coefficients,
then passed to inverse limits and rationalized.}
As in Subsection~\ref{subsec:frobenius-neighbors}, crystalline Frobenius is
the pullback of the special fiber's \(q\)-power endomorphism: Witt
Frobenius on \(W(\Fq)=\Z_q\) is the identity
\cite[Tags~07N3 and~07N5]{StacksProject}.
Crystalline representations give weakly admissible filtered modules
\cite[Theorem~9.3.4]{BrinonConrad2009}. The degree-one Hodge filtration
has only indices \(0\) and \(1\), and functoriality gives the assertion
about extending endomorphisms.

The extension by \(L\) changes coefficients, not the geometric base
field. With Frobenius acting trivially on this coefficient factor, the
underlying \(\Q_q\)-filtered Frobenius module
\[
\bigl(D_A\otimes_{\Q_q}L,\ \varphi_A\otimes1,\
       F^iD_A\otimes_{\Q_q}L\bigr)
\]
is a direct sum of \([L:\Q_q]\) copies of \(D_A\), and is weakly
admissible. Since \(p\) and \(1-p\) both preserve the filtration and
commute with Frobenius, their images split it as filtered Frobenius
modules. Weak admissibility gives \(t_H\leq t_N\) on each image.
The Hodge and Newton numbers are additive, and are equal for the whole
module. They are therefore equal on each image separately.

Let \(m=\dim_{\Q_q}p(D_A\otimes L)\). A pure slope \(t\) gives
\(t_N=tm\) after restriction of scalars: all \(\Q_q\)-conjugates of a
\(q\)-adic algebraic number have the same valuation, including when
\(L/\Q_q\) is ramified. On the other hand,
\[
t_H=\dim_{\Q_q}F^1\bigl(p(D_A\otimes L)\bigr),
\]
because the only filtration indices are \(0\) and \(1\). Equality of
these two numbers proves the assertion. The general coefficient
formalism is also described in
\cite[Proposition~2.17(ii) and Corollary~2.22]{Hellmann2011}.
\end{proof}

The requirement on the idempotent is essential. For example,
\(\varphi=\diag(1,q)\) on \(\Q_qe_0\oplus\Q_qe_1\), with only
filtration indices \(0,1\) and \(F^1=\langle e_0+e_1\rangle\), is
weakly admissible, but its slope-one line has induced \(F^1=0\).
Its slope projector does not preserve the filtration. The Hecke
projectors do preserve it, because they are formed from extending
algebraic endomorphisms.

We can therefore obtain a pure Hodge type by proving that a full
generalized Hecke space has slope zero or one. The finite-order central
relation controls the common scalar in a full extension of the active
projective parameter, and the Frobenius polynomial then determines the
slope.

Fix an active character \(\eta\) with projective parameter
\[
\{q^{1/2},q^{-1/2},e'\},
\]
where \(e'=e\) is supplied by Proposition~\ref{prop:theta-satake}, or
\(e'=\pm\sigma e\) is the parameter of a scalar-conjugate character.
Thus \(e'\) is algebraic, all its complex conjugates have modulus one,
and \(\nu_\iota(e')\in\{-3/2,3/2\}\). For a full character \(\chi\)
extending \(\eta\), its active \(\GL_3\) parameter is a scalar
multiple of this triple. Absorb the numerical value of \(C_0\) into
that scalar and write the resulting triple as
\begin{equation}\label{eq:full-satake-triple}
(\alpha_1,\alpha_2,\alpha_3)
=b(q^{1/2},q^{-1/2},e').
\end{equation}
The determinant coordinate is the character value of
\(D_3=C_0^3T_3\). Lemma~\ref{lem:central-finite-order} therefore gives
\begin{equation}\label{eq:satake-central-relation}
b^3e' \text{ is a root of unity}.
\end{equation}
In particular \(b\) is algebraic and every complex conjugate of \(b\)
has modulus one.

Write \(e_i\) for the elementary symmetric polynomials in three
variables. The raw minuscule double cosets, with the integral subgroup
of volume one, have Satake transforms
\(q e_1,\ q e_2,\ e_3\); this is
\cite[Equation~(3.14) and \S6]{Gross1998} for \(\GL_3\).
Since \(D_i=C_0^iT_i\), absorption of the \(C_0\)-value into \(b\)
gives
\begin{equation}\label{eq:raw-D-character}
\chi(D_1)=q e_1(\alpha),\qquad
\chi(D_2)=q e_2(\alpha),\qquad
\chi(D_3)=e_3(\alpha).
\end{equation}
These are numerical character values. On \(H_\chi\) the operators
may still have nilpotent parts.

\begin{proposition}\label{prop:active-hodge-types}
Let \(\eta\) be an active numerical character with projective parameter
\(\{q^{1/2},q^{-1/2},e'\}\) supplied by
Proposition~\ref{prop:theta-satake} or by scalar conjugation, and put
\[
t_\eta=\frac12-\frac13\nu_\iota(e')\in\{0,1\}.
\]
For every full numerical character \(\chi\) extending \(\eta\), the
space \(H_\chi\otimes_{\Qbar}\C\) is contained in \(H^{1,0}(A_S)\)
if \(t_\eta=1\), and in \(H^{0,1}(A_S)\) if \(t_\eta=0\).
Consequently \(H^{\mathrm{act}}_\eta\otimes_{\Qbar}\C\) has the same
pure Hodge type.
\end{proposition}

\begin{proof}
Fix a full character \(\chi\) extending \(\eta\); if there is none,
the assertion is vacuous.
We first use Weil weight to select the Frobenius root. We then apply
Lemma~\ref{lem:filtered-hecke-summands} and compare the resulting
filtration ranks with those on the complex Betti realization.

At the fixed completion \(K_v=\Q_q\), write
\[
D_S=H^1_{\mathrm{dR}}((A_S)_{K_v}/\Q_q)
\]
and let \(\varphi\) be its crystalline Frobenius. If \(\Phi\) is the
\(q\)-power Frobenius endomorphism of \(\widetilde A_S/\Fq\) used in
Proposition~\ref{prop:albanese-congruence}, then
\(\varphi=\Phi^*\) under crystalline comparison. The Hecke
endomorphisms extend over \(\Z_q\); their reductions commute with
\(\Phi\), so their pullbacks commute with \(\varphi\).

By additivity of pullback on abelian degree-one cohomology, the integral identity of
Proposition~\ref{prop:albanese-congruence} gives
\begin{equation}\label{eq:cohomological-frobenius-polynomial}
\varphi^3-D_1^*\varphi^2+qD_2^*\varphi-q^3D_3^*=0.
\end{equation}
Pullback reverses composition, but all operators in this formula
commute.

Choose a prime \(\ell\ne q\) and put
\(V_\ell\widetilde A_S=T_\ell\widetilde A_S\otimes_{\Z_\ell}\Q_\ell\).
Weil's theorem for abelian varieties over \(\Fq\) gives
\[
 P_\Phi(X)=\det\bigl(X-\Phi\mid V_\ell\widetilde A_S\bigr)\in\Z[X],
\]
all of whose complex roots have absolute value \(q^{1/2}\).
The prime-to-\(q\) Tate-module action on endomorphisms is faithful.
Cayley--Hamilton therefore gives \(P_\Phi(\Phi)=0\) in the
endomorphism ring, and degree-one additivity gives
\(P_\Phi(\varphi)=0\).  These statements hold for an arbitrary
abelian variety over \(\Fq\); see
\cite[Lemma~12.2, Proposition~12.9, and Theorem~19.1(c) with its proof]{Milne1986AbelianVarieties}.
No CM or ordinary hypothesis on \(\widetilde A_S\) is involved.
This uses the \(q\)-power endomorphism and its pullback, not a
reciprocal Galois-Frobenius convention.

Transport the algebraic coefficients of \(p_\chi\) and \(\chi\) by
\(\iota\). On the corresponding generalized summand of
\(D_S\otimes_{\Q_q}\overline{\Q}_q\), the commuting Hecke maps and
\(\varphi\) can be simultaneously triangularized. Equations
\eqref{eq:cohomological-frobenius-polynomial}
and~\eqref{eq:raw-D-character} imply that every diagonal Frobenius
eigenvalue is a root of
\[
X^3-qe_1(\alpha)X^2+q^2e_2(\alpha)X-q^3e_3(\alpha)
=\prod_{j=1}^3(X-q\alpha_j).
\]
Thus the candidates, before applying \(\iota\), are
\begin{equation}\label{eq:frobenius-candidates}
bq^{3/2},\qquad bq^{1/2},\qquad qbe'.
\end{equation}
Each is algebraic. If its image under \(\iota\) is a root of the
rational polynomial \(P_\Phi\), the algebraic number itself is a
root of \(P_\Phi\). We may therefore apply the complex
absolute-value condition. The three absolute values are respectively
\[
q^{3/2},\qquad q^{1/2},\qquad q,
\]
because every complex conjugate of \(b\) and \(e'\) has modulus one.
Only \(bq^{1/2}\) remains.

The weight condition has now selected one Newton slope on the entire
generalized summand. Equation~\eqref{eq:satake-central-relation} gives
\begin{equation}\label{eq:active-hodge-slope}
\nu_\iota(bq^{1/2})
=\frac12-\frac13\nu_\iota(e')=t_\eta.
\end{equation}
This conclusion concerns every Frobenius eigenvalue; Jordan blocks do
not change it. Choose a finite extension \(L/\Q_q\) containing the
coefficients of the transported idempotent and numerical character.
The idempotent \(p_\chi\)
is a polynomial in extending Hecke maps, so it preserves the Hodge
filtration and commutes with \(\varphi\). Lemma
\ref{lem:filtered-hecke-summands} makes its filtration pure of index
\(t_\eta\).

It remains to relate this filtered rank to the Hodge type of the same
complex generalized space. Put
\[
V_{\mathrm{dR}}=H^1_{\mathrm{dR}}(A_S/K).
\]
Its filtration and Hecke operators are defined over the fixed field
\(K\). For a numerical character \(\chi\), choose a polynomial
expression for \(p_\chi\), and let \(k_\chi\subset\Qbar\) be the
actual compositum of \(K\) and its algebraic coefficients. Evaluating
that polynomial in the fixed operators gives an idempotent on
\(V_{\mathrm{dR}}\otimes_K k_\chi\). Its induced filtration is its
projected filtration, since the idempotent preserves \(F^\bullet\).
The dimensions of its image and of this filtration are ranks of
matrices over \(k_\chi\);
they are unchanged under the restriction of the fixed embedding
\(\iota\) to \(k_\chi\).

Over the given embedding \(K\subset\C\), Betti--de Rham comparison
intertwines the same operators and polynomial. It identifies these
ranks with the filtration ranks on \(H_\chi\otimes_{\Qbar}\C\).
For a scalar-conjugate numerical character \(\sigma\chi\), repeat
this argument with \(p_{\sigma\chi}\) in the same \(K\)-defined
operators. On Betti cohomology their rational matrices give
\(\sigma H_\chi=H_{\sigma\chi}\). We have therefore computed the
type of the scalar-conjugated Betti space on the same complex variety.
Each conjugate projector is evaluated separately in the fixed
\(K\)-defined operators. This works even when \(K\) intersects the
character field, since it uses ranks over their actual compositum.

Finally the complex Hecke operators preserve both \(H^{1,0}\) and
\(H^{0,1}\). Hence a full projected \(F^1\) means that the whole
generalized space is holomorphic, and zero projected \(F^1\) means
that it is antiholomorphic. Equation~\eqref{eq:active-hodge-slope}
depends only on \(\eta\), so summing all full extensions gives the
assertion for \(H^{\mathrm{act}}_\eta\).
\end{proof}

\subsection{Rational CM sources for the theta classes}\label{subsec:rational-cm-sources}

The preceding proposition determines the type of every vector in each
scalar-conjugate active Hecke space. We can therefore apply scalar
descent to maps with arbitrary algebraic Betti vectors in these spaces.
This will prove the rational-map span asserted in
Theorem~\ref{thm:cm-extraction}.

\begin{proof}[Proof of Theorem~\ref{thm:cm-extraction}]
Use the Albanese identification
\(H^1(A_S,\Q)\simeq H^1(S,\Q)\). Extend a theta class from \(S_0\)
by zero on the other components of \(S\). This is compatible with
the active algebra, because Lemma~\ref{lem:active-hecke-theta}
makes that algebra component-preserving. The class is holomorphic
by Proposition~\ref{prop:theta-one}. Decompose it into the active
generalized character spaces of
Equation~\eqref{eq:active-generalized-space}. For each nonzero component,
let \(\eta\) be its active character and \(e\) its parameter from
Proposition~\ref{prop:theta-satake}. For every
\(\sigma\in\Gal(\Qbar/\Q)\), put \(g=\sigma|_E\). That proposition
and Proposition~\ref{prop:active-hodge-types} give
\begin{equation}\label{eq:cm-sign-to-hodge-type}
\nu_\iota(\sigma e)=-\frac32v(g),\qquad
t_{\sigma\eta}
=\frac12-\frac13\nu_\iota(\sigma e)
=\frac{1+v(g)}2.
\end{equation}
The possible sign in the conjugate parameter \(e'=\pm\sigma e\)
has valuation zero. Consequently the entire active space for
\(\sigma\eta\) is of type \((1,0)\) when \(v(g)=1\) and of type
\((0,1)\) when \(v(g)=-1\). This applies to all of its vectors,
not just to theta vectors. At the identity it gives type \((1,0)\),
in agreement with \(v(1)=1\) and the holomorphy of the theta class.

Each active space is defined over \(\Qbar\) in the rational Betti
structure, because its idempotent has algebraic coefficients.
Choose \(0\ne x\in H^{\mathrm{act}}_\eta\) and
\(0\ne u\in U(v)_1(\Qbar)\). Define a \(\Qbar\)-linear map
\[
h:U(v)_{\Qbar}\longrightarrow H^1(S,\Qbar)
\]
by \(h(u)=x\) and by setting it equal to zero on every other
embedding line of \(U(v)_{\Qbar}\). This is an element of
\[
\Hom_{\Q}\bigl(U(v),H^1(S,\Q)\bigr)\otimes_{\Q}\Qbar.
\]
Scalar conjugation by \(\sigma\) moves its nonzero source line to
\(U(v)_g\) and its target vector to
\(\sigma x\in H^{\mathrm{act}}_{\sigma\eta}\). The source line has
the type prescribed by \(v(g)\), and
Equation~\eqref{eq:cm-sign-to-hodge-type} gives the same type for
the target. All other source lines map to zero. Hence every scalar
conjugate of \(h\) is a Hodge map.

Lemma~\ref{lem:trace-descent}, applied to the complex-linear
subspace of Hodge maps, expresses \(h\) as a \(\Qbar\)-linear
combination of rational Hodge maps. Algebraic-Betti vectors span
each complex active space, so taking complex spans covers every
active component of the original theta class, including classes
with transcendental coefficients. Restriction from \(S\) to
\(S_0\) is a rational Hodge map and gives
Equation~\eqref{eq:cm-extraction-span}.

For a finite collection of inputs, all level choices are made
before the auxiliary prime is chosen, and only finitely many primes
are excluded. The preceding argument therefore applies to that
collection at one sufficiently deep level. If
\(\pi:S'_0\to S_0\) is a deeper principal cover, then
\(\pi^*\) is a rational Hodge map and the theta classes pull back
to their classes on \(S'_0\). Applying \(\pi^*\) to
Equation~\eqref{eq:cm-extraction-span} proves persistence.
\end{proof}

By Lemma~\ref{lem:weight-one-maps}, each rational map in
Equation~\eqref{eq:cm-extraction-span} is, up to a rational multiple,
pullback along a morphism \(S_0\to B(v)\).

\begin{corollary}[Finite-family CM extraction]\label{cor:finite-family-extraction}
Fix \(E,V,D,\Lambda\) and the principal lattice levels of
Section~\ref{sec:arithmetic}.
Let \(v_1,\ldots,v_m:G_E\to\{1,-1\}\) be odd with \(v_j(1)=1\).
For each \(j\), choose a finite family of finite adelic Schwartz inputs
in the multiplicity-line construction
\(d_j(\tau)=v_j(\tau^{-1})\).  Let \(X_{\mathrm{base}}\) be a
principal quotient at a torsion-free level in this framework.
There is a common deeper principal cover \(Y\to X_{\mathrm{base}}\),
with \(Y\) a smooth connected projective surface, on which every
associated theta one-form class from the \(j\)-th family belongs to
\[
\operatorname{span}_{\C}\left\{
f\bigl(U(v_j)_1\bigr):
f\in\Hom_{\mathrm{Hdg}}\bigl(U(v_j),H^1(Y,\Q)\bigr)
\right\}.
\]
The maps are rational Hodge maps, with their images viewed after
extension to \(\C\).  The conclusion persists after pullback to every
deeper principal level.  The auxiliary primes used for different
\(v_j\) may differ.
\end{corollary}

\begin{proof}
Apply Theorem~\ref{thm:cm-extraction} to each finite input family.
Choose a principal level contained in the resulting finitely many
principal levels and in the level of \(X_{\mathrm{base}}\), and deepen
it until Lemma~\ref{lem:reference-moduli} applies.  Its quotient \(Y\)
is smooth, connected, and projective.  The reference moduli construction
depends on \(V\), not on the multiplicity line, so it applies to this
common level.  The persistence in Theorem~\ref{thm:cm-extraction}
pulls all the rational-map spans to \(Y\) and then to every deeper
principal level.  Each application chooses its auxiliary prime only
after its own finite family and level have been fixed; the primes need
not agree.
\end{proof}
